% !TeX root = courtade-kumar.tex
% !TEX TS-program = pdfLaTeX
\section{Introduction}
\label{sec:introduction}
Let $X$ be distributed uniformly on $\cube$ and let $Y$
be obtained by copying $X$ and then flipping each of its
coordinates $Y_i$ with probability $p$, independently
for $i\in[n]$.
In \cite{KumarC2013}, Courtade and Kumar ask, what is the
best strategy to output a bit $B$, having access to $Y$
only, so that $\muti{B}{X}$, the information revealed about
$X$, is maximized.

They conjecture that taking $B=Y_i$ for some $i\in[n]$
maximizes the information conveyed about $X$. In particular,
their Conjecture 1 states that 
\begin{equation}
\label{eq:ck}\tag{(CK)}
\muti{B}{X}\le\muti{Y_1}{X} = 1 - \entb{p}.
\end{equation}
They further remark that, despite its resemblance to a
freebie homework exercise, the statement appears to
involve substantial depth. Indeed, despite extensive
efforts and partial progress
\cite{AnantharamBCJN2017,Arikan2009,BarnesO2020,ChenN2024,
FischerKRSS2004,JavanmardW2026,KindlerOW2015,LiM2021,
OrdentlichSW2015,ParnasRS2002,PichlerMP2016,Samorodnitsky2015,
Samorodnitsky2016,Sarbu2017,Sarbu2017b,YangW2018,Yu2024},
a full resolution of the conjecture remained elusive.

In this work, we show that if we study the inequality
parametrized by $\mu=\E\sparen{B}$, a delicate refinement
of \eqref{eq:ck} emerges, which admits an induction on $n$,
the number of coordinates of $X,Y$. In particular, we prove:

\begin{theorem}
\label{thm:ck+}
Let $X$ be a uniform random bit string and $Y$ be obtained
by copying $X$ and then keeping each coordinate as-is with
probability $\rho$, replacing it with a uniform random bit
with probability $1-\rho$, independently.
Let $B$ be a random bit obtained as a function of $Y$ with
expectation $\mu$. Writing $\nu = 4\mu(1-\mu)$, we have
\begin{equation*}
\muti{B}{X}\le \rho^2\nparen{\entb{\mu}-\nu}
  + \nu\muti{Y_1}{X}
\end{equation*}
for any $\rho,\mu\in[0,1]$.
\end{theorem}
Two remarks are in order. First, note that an alternative
description of $Y$ is to keep each coordinate of $X$ with
probability $\rho = 1-2p$ and replace it with a uniform random
bit with probability $1-\rho$. Flipping a bit with
probability $p$ or re-randomizing with probability $1-\rho$
leads to the same distribution. Secondly, setting
\newcommand{\Q}{Q_\rho(\mu)}
\begin{equation}
\label{eq:q}\tag{(Q)}
\Q = \rho^2\nparen{\entb{\mu}-\nu}
  + \nu\muti{Y_1}{X},
\end{equation}
we have $\Q\le \muti{Y_1}{X}$ for all $\rho,\mu\in[0,1]$,
whose proof we defer to \autoref{sec:proof}.

\subsection{The continuous case}
We note that $\muti{B}{X}\le \Q$ is a surprisingly sharp
bound, which is already violated by the most natural
generalization of the problem.
Let $X_\circ$ be distributed uniformly on the unit cube
$[0,1]^n$ and $Y_\circ$ be obtained by copying the
corresponding coordinate of $X_\circ$ with probability
$\rho$ and drawing uniformly from $[0,1]$ with probability
$1-\rho$, for each coordinate $i\in[n]$ independently.
Let $B_\circ$ be a random bit obtained as a function of
$Y_\circ$ with $\mu=\E\sparen{B_\circ}$.

Denote by $C_\mu$ the indicator of the event
$Y_\circ[1]\le \mu$, that is, the first coordinate of
$Y_\circ$ being at most $\mu$. Note that
$\E\sparen{C_\mu}=\mu$. In \autoref{lem:ck+vscontinuous} we
show that $\Q < \muti{C_\mu}{X_\circ}$ for $0<\mu<1$,
$\mu\ne1/2$, and $0<p<1/2$, which means the threshold
functions already exceed the bound $\Q$. We further
conjecture the following.
\begin{conjecture}
\label{conj:continuous}
Let $X_\circ$, $Y_\circ$ and $B_\circ$ be as constructed
and $C_\mu$ be the indicator of the event $Y_\circ[1]\le\mu$
as above. It holds that
\begin{equation*}
\muti{B_\circ}{X_\circ}\le\muti{C_\mu}{X_\circ}
\end{equation*}
for all $\rho, \mu\in [0,1]$.
\end{conjecture}

Note that a proof of \autoref{conj:continuous} directly
implies \eqref{eq:ck}---though not our
\autoref{thm:ck+}---since given a strategy to output $B$
from $Y$ revealing $\muti{X}{B}$ bits of information, we can
quantize $Y_\circ$ by rounding each coordinate to $0$ or $1$
with threshold $1/2$ and apply the strategy to obtain
$B_\circ$. We get $\muti{B_\circ}{X_\circ}=\muti{B}{X}$.

\subsection{Notation}
\label{sec:notation}
All logarithms are base $2$, except $\ln$, which denotes the
natural logarithm. We write
\begin{align*}
 \entb{a}     &=a\log \frac{1}{a}+ (1-a)\log\frac{1}{1-a},\\
 \kldivb{a}{b}&=a\log\frac ab+(1-a)\log\frac{1-a}{1-b}
\end{align*}
for the binary entropy and divergence, respectively, with endpoint
values understood by continuity.
For $a\in[0,1]$ and $b\in[1/2,1)$, we repeatedly use
\begin{equation}
\tag{(U)}
\label{eq:U}
 \frac{2(a-b)^2}{\ln2}
 \le\kldivb{a}{b}
 \le\frac{\ln(1/b)}{(1-b)^2\ln2}(a-b)^2.
\end{equation}
The identity $\kldivb{a}{b}=\kldivb{1-a}{1-b}$ covers
$0<b<1/2$ as well. The lower bound also holds for $b=0,1$.

For the lower bound, the second derivative of the divergence in
$a$ is $1/(a(1-a)\ln2)\ge4/\ln2$, and its value and first
derivative vanish at $a=b$.
For the upper bound, Taylor's integral formula gives
\[
 \frac{\kldivb{a}{b}}{(a-b)^2}
 =\frac1{\ln2}\int_0^1
   \frac{1-s}{\nparen{b+s(a-b)}\nparen{1-b-s(a-b)}}
   \,\mathrm{d}s.
\]
The quotient extends continuously to $a=b$ and is convex in $a$,
since $1/(x(1-x))=1/x+1/(1-x)$ is convex on $(0,1)$.
It therefore attains its maximum at an endpoint.
The value at $a=1$ is at least the value at $a=0$ because
\[
 -b^2\ln b+(1-b)^2\ln(1-b)\ge0.
\]
Indeed, the expression vanishes at $b=1/2$ and $b=1$, and its
second derivative is $2\ln((1-b)/b)\le0$ on this interval.
The value at $a=1$ is the coefficient in the upper bound of
\eqref{eq:U}. Strict convexity of the quotient also makes this
bound strict when $0<a<1$ and $a\ne b$.

We will also use the Taylor expansion
\begin{equation}
\label{eq:binary-divergence-series}
 \kldivb{\frac{1+u}{2}}{1/2}
 =\frac1{2\ln2}\sum_{k\ge1}\frac{u^{2k}}{k(2k-1)}
 \qquad(-1\le u\le1).
\end{equation}
Its value and first derivative vanish at $u=0$, and its second
derivative is $1/((1-u^2)\ln2)$.
Expanding this derivative as a geometric series and integrating
twice proves the identity for $|u|<1$; absolute convergence and
continuity include the endpoints.
